\section{September 2, 2026}

\subsection{Linear transformations determined by a basis}

The following result shows that prescribing the values of a linear
transformation on a basis completely determines the transformation.

\begin{lemma}[Linear map lemma]\label{lem:linear-map}
  Suppose that $\{v_1,\ldots,v_n\}$ is a basis for $V$, and let
  $w_1,\ldots,w_n$ be arbitrary vectors in $W$. Then there exists a
  unique linear transformation $T\colon V\to W$ such that
  \[
    T(v_k)=w_k
  \]
  for each $k=1,\ldots,n$.
\end{lemma}

\begin{proof}
  Every $v\in V$ has a unique representation
  \[
    v=\sum_{i=1}^n c_iv_i.
  \]
  Define
  \[
    T(v)=\sum_{i=1}^n c_iw_i.
  \]
  The uniqueness of the coordinates $c_1,\ldots,c_n$ shows that this
  definition is well-defined. In particular, taking $v=v_k$ gives
  $T(v_k)=w_k$.

  To prove that $T$ is linear, write
  \[
    u=\sum_{i=1}^n a_iv_i
    \qquad\text{and}\qquad
    v=\sum_{i=1}^n c_iv_i.
  \]
  Then
  \begin{align*}
    T(u+v)
    &=T\left(\sum_{i=1}^n(a_i+c_i)v_i\right) \\
    &=\sum_{i=1}^n(a_i+c_i)w_i \\
    &=\sum_{i=1}^n a_iw_i+\sum_{i=1}^n c_iw_i \\
    &=T(u)+T(v).
  \end{align*}
  Similarly, for every $\alpha\in\bb{F}$,
  \begin{align*}
    T(\alpha v)
    &=T\left(\sum_{i=1}^n\alpha c_iv_i\right) \\
    &=\sum_{i=1}^n\alpha c_iw_i
      =\alpha\sum_{i=1}^n c_iw_i
      =\alpha T(v).
  \end{align*}
  Thus $T$ is linear.

  It remains to prove uniqueness. Suppose that $T'\colon V\to W$ is
  another linear transformation satisfying $T'(v_k)=w_k$ for every
  $k$. If $v=\sum_{i=1}^n c_iv_i$, then
  \[
    T'(v)
    =\sum_{i=1}^n c_iT'(v_i)
    =\sum_{i=1}^n c_iw_i
    =T(v).
  \]
  Hence $T'=T$.
\end{proof}

\subsection{Matrix representations of linear transformations}

An \emph{ordered basis} is a basis whose vectors are listed in a fixed
order. Let
\[
  \mathcal B=(v_1,\ldots,v_n)
  \qquad\text{and}\qquad
  \mathcal C=(w_1,\ldots,w_m)
\]
be ordered bases for $V$ and $W$, respectively. For
$v=x_1v_1+\cdots+x_nv_n$, write
\[
  [v]_{\mathcal B}
  =\begin{bmatrix}x_1\\ \vdots\\ x_n\end{bmatrix}
\]
for the coordinate vector of $v$ with respect to $\mathcal B$.

\begin{theorem}[Matrix representation theorem]
  Let $V$ and $W$ be vector spaces over $\bb{F}$ with ordered bases
  $\mathcal B=(v_1,\ldots,v_n)$ and
  $\mathcal C=(w_1,\ldots,w_m)$. For every linear transformation
  $T\colon V\to W$, there is a unique matrix
  \[
    A=[A_{ij}]
    =\begin{bmatrix}
      A_{11} & \cdots & A_{1n} \\
      \vdots &        & \vdots \\
      A_{m1} & \cdots & A_{mn}
    \end{bmatrix}
    \in M_{m\times n}(\bb{F})
  \]
  such that
  \[
    T(v_j)=\sum_{i=1}^m A_{ij}w_i
  \]
  for each $j=1,\ldots,n$. Conversely, every matrix
  $A\in M_{m\times n}(\bb{F})$ determines a unique linear transformation
  $T\colon V\to W$ through this formula.
\end{theorem}

\begin{proof}
  Let $T\colon V\to W$ be linear. Because $\mathcal C$ is a basis for
  $W$, each vector $T(v_j)$ has a unique representation
  \[
    T(v_j)=A_{1j}w_1+\cdots+A_{mj}w_m
           =\sum_{i=1}^m A_{ij}w_i.
  \]
  Use these coordinate vectors as the columns of $A$:
  \[
    A
    =\begin{bmatrix}
      \vert &        & \vert \\
      [T(v_1)]_{\mathcal C} & \cdots & [T(v_n)]_{\mathcal C} \\
      \vert &        & \vert
    \end{bmatrix}.
  \]
  This constructs the required matrix. If $B=[B_{ij}]$ were another
  matrix with the same property, then for every $j$,
  \[
    \sum_{i=1}^m B_{ij}w_i
    =T(v_j)
    =\sum_{i=1}^m A_{ij}w_i.
  \]
  The uniqueness of coordinates with respect to $\mathcal C$ gives
  $B_{ij}=A_{ij}$ for every $i$ and $j$. Hence $B=A$.

  Conversely, let $A=[A_{ij}]$ be an arbitrary $m\times n$ matrix and
  define
  \[
    \widetilde w_j=\sum_{i=1}^m A_{ij}w_i
    \qquad (j=1,\ldots,n).
  \]
  By the linear map lemma, there is a unique linear transformation
  $T\colon V\to W$ such that $T(v_j)=\widetilde w_j$ for every $j$.
  Thus $A$ determines a unique linear transformation with the required
  property.
\end{proof}

If
\[
  v=x_1v_1+\cdots+x_nv_n,
\]
then linearity gives
\begin{align*}
  T(v)
  &=\sum_{j=1}^n x_jT(v_j) \\
  &=\sum_{j=1}^n x_j\sum_{i=1}^m A_{ij}w_i \\
  &=\sum_{i=1}^m\left(\sum_{j=1}^n A_{ij}x_j\right)w_i.
\end{align*}
Consequently,
\[
  [T(v)]_{\mathcal C}
  =A[v]_{\mathcal B}.
\]
Thus multiplication by the matrix $A$ computes the coordinates of the
image of a vector from the coordinates of the original vector.
